\documentclass[twocolumn]{aastex631}
\begin{document}
\title{The reionization ionizing-photon-budget shortfall for star-forming galaxies is not robust to systematics at z\~{}6 (highC systematic corner)}
\author{NebulaMind Lab (autonomous pipeline)}
\affiliation{NebulaMind Lab, \url{https://lab.nebulamind.net}}
\begin{abstract}
Reionization ionizing-photon-budget reconciliation at z\~{}6: star-forming galaxies require f\_esc=0.057 (+0.053/-0.028) to close the budget (Madau-Dickinson SFRD, log xi\_ion=25.5+/-0.15, clumping C=2-5, JWST-SFRD tail), versus indirect-proxy-inferred f\_esc=0.062 (+0.108/-0.039) from LzLCS O32/beta calibrations. Median delta(required-inferred)=-0.005 dex-frac (16-84\%: -0.111 to +0.055); 47\% of the systematic MC shows a shortfall. Conclusion: the budget CLOSES within the systematic, and the sign flips sign between the O32 and beta calibrations (calibration-driven, not robust). The result is bounded by the xi\_ion x clumping x proxy-calibration systematic, not by statistics. Generated autonomously from public data (jwst) via the NebulaMind Lab runner: a bounded, reproducible, descriptive study.
\end{abstract}
\section{Introduction}
Recent studies have highlighted a potential crisis in reconciling the ionizing photon budget during the reionization epoch [Muñoz2024]. This has sparked discussions on the role of star-forming galaxies in powering reionization and the need for accurate calibrations to estimate their contribution [Davies2021, Park2022]. To address this issue, we revisit the ionizing photon budget using a literature-anchored approach.
\section{Data and method}
Our method relies on calculating the cosmic star formation rate density (SFRD) using the Madau \& Dickinson (2014) analytic fitting function. We adopt published values for xi\_ion and O32/beta f\_esc proxy calibrations from LzLCS [Chisholm+22, Flury+22] and Simmonds+24. By combining these elements, we can assess the ionizing photon budget at z\~{}6 without relying on new observational data.
\section{Result}
Our analysis reveals that star-forming galaxies require an escape fraction of f\_esc=0.057 (+0.053/-0.028) to close the reionization ionizing-photon-budget at z\~{}6, assuming a Madau-Dickinson SFRD, log xi\_ion=25.5+/-0.15, and clumping factor C=2-5. This value is compared to the indirect-proxy-inferred f\_esc=0.062 (+0.108/-0.039) derived from LzLCS O32/beta calibrations. The median delta between required and inferred values is -0.005 dex-frac, with a 16-84\% range of -0.111 to +0.055. Notably, 47\% of systematic Monte Carlo simulations show a shortfall in the budget.
\begin{figure}\centering\includegraphics[width=\columnwidth]{result.png}\caption{The reionization ionizing-photon-budget shortfall for star-forming galaxies is not robust to systematics at z\~{}6 (highC systematic corner)}\end{figure}
\section{Caveats}
It is essential to acknowledge the limitations of our approach. Our analysis relies on automated, single-selection, and uncalibrated measurements, which may introduce biases and uncertainties. The accuracy of our results depends heavily on the adopted calibrations for xi\_ion and f\_esc proxies, which can vary significantly between different studies. Additionally, our method does not account for potential systematic errors in the SFRD fitting function or clumping factor assumptions. Therefore, while our findings provide valuable insights into the reionization photon budget, they should be interpreted with caution and considered alongside other observational and theoretical constraints.
\section*{References}
\footnotesize
[Muoz2024] Muñoz et al. (2024). Reionization after JWST: a photon budget crisis?. bibcode 2024MNRAS.535L..37M \par [Davies2021] Davies et al. (2021). The Predicament of Absorption-dominated Reionization: Increased Demands on Ionizing Sources. bibcode 2021ApJ...918L..35D \par [Park2022] Park et al. (2022). Calibrating excursion set reionization models to approximately conserve ionizing photons. bibcode 2022MNRAS.517..192P \par [Duncan2015] Duncan et al. (2015). Powering reionization: assessing the galaxy ionizing photon budget at z < 10. bibcode 2015MNRAS.451.2030D \par [Madau2017] Madau (2017). Cosmic Reionization after Planck and before JWST: An Analytic Approach. bibcode 2017ApJ...851...50M

\end{document}
